% Source: https://artofproblemsolving.com/community/c3h2425117_100th_post
% Stable identity: aops-topic-2425117-q01
For each point of the coating, look at the point of the box nearest to it. It lies in the interior of a face, on an edge, or at a corner, and this splits the coating into pieces:
\begin{itemize}
\item over each face, a slab of thickness $2$ whose base is the face;
\item along each edge, a quarter of a cylinder of radius $2$ whose length is the edge;
\item at each corner, an eighth of a ball of radius $2$.
\end{itemize}
The cross-section below, perpendicular to the edges of length $3$, shows the slabs and the quarter discs (quarter cylinders) between them.

\begin{asy}
usepackage("amsmath");
size(5.5cm);
defaultpen(fontsize(10pt));
real t=2, w=2, h=1;
path outer=(0,-t)--(w,-t)..arc((w,0),t,-90,0)--(w+t,h)..arc((w,h),t,0,90)--(0,h+t)..arc((0,h),t,90,180)--(-t,0)..arc((0,0),t,180,270)--cycle;
fill(outer,rgb(0.8,0.85,0.95));
filldraw((0,0)--(w,0)--(w,h)--(0,h)--cycle,white,linewidth(0.9));
draw(outer,linewidth(0.8));
draw((0,0)--(0,-t),dashed); draw((w,0)--(w,-t),dashed); draw((w,0)--(w+t,0),dashed); draw((w,h)--(w+t,h),dashed);
draw((w,h)--(w,h+t),dashed); draw((0,h)--(0,h+t),dashed); draw((0,h)--(-t,h),dashed); draw((0,0)--(-t,0),dashed);
draw((w/2,0)--(w/2,-t),Arrows(4)); label("$2$",(w/2,-t/2),E);
draw((w,h)--(w,h)+t*dir(45),Arrows(4)); label("$2$",(w,h)+t/2*dir(45),NW);
label("box",(w/2,h/2)); label("slab",(w/2,h+t/2)); label("quarter",(w+1.25,-1.15),fontsize(8pt)); label("disc",(w+1.25,-1.55),fontsize(8pt));
\end{asy}

\textbf{Slabs.} The surface area of the box is $2(1\cdot2+2\cdot3+1\cdot3)=22$, so the slabs have volume $22\cdot2=44$.

\textbf{Quarter cylinders.} There are $4$ edges of each length $1$, $2$ and $3$, of total length $4(1+2+3)=24$. Four quarter cylinders of the same length would make a whole cylinder, so these have volume $\dfrac14\pi\cdot2^2\cdot24=24\pi$.

\textbf{Corners.} The $8$ eighth-balls make one whole ball of radius $2$, of volume $\dfrac43\pi\cdot2^3=\dfrac{32\pi}{3}$.

The volume of titanium is
\[44+24\pi+\dfrac{32\pi}{3}=\boxed{44+\dfrac{104\pi}{3}\text{ cubic inches}}\approx152.9\text{ cubic inches}.\]
